\subsection*{Introduction}

The neutronics class (see {\itshape \hyperlink{neutronics_8H}{neutronics.\+H}}) is a high-\/level class that contains essential data and variables that are common to various neutronics models. In particular, the neutronics class handles the variables that are included in the {\itshape constant/neutro\+Region/reactor\+State} dictionary.

 {\bfseries The {\itshape reactor\+State} dictionary}

The {\itshape reactor\+State} dictionary is found under the {\itshape time\+Folder/uniform/} sub-\/folder. It included essentially 4 keywords\+: 
\begin{DoxyItemize}
\item {\itshape keff} is used in the spatial kinetics solvers as initial guess for keff when doing an eigenvalue calculation. It is then updated automatically at each time step (i.\+e., at each power iteration) with the calculated value of keff. When performing a transient calculation with the spatial kinetics solvers, {\itshape keff} is instead used to divide the neutron source term and is not updated during the simulation. Typically, to run spatial kinetics transient simulations, one first runs an eigenvalue calculation. The resuling {\itshape keff} will be the one that makes the reactor critical in a subsequent transient simulation. {\itshape keff} is diregarded by the point kinetics sub-\/solver. 
\item {\itshape p\+Target} is used in the spatial kinetics solvers as target power when doing an eigenvalue calculation. It is also used by the point kinetics sub-\/solver, but only to correctly plot results. As power, Ge\+N-\/\+Foam uses what it finds under power\+Density of the neutro\+Region, or under the power\+Density of the fluid\+Region if it does not find a power\+Density in the neutro\+Region. To correctly plot point kinetics result, p\+Target must be consistent with the mentioned power densities. 
\item {\itshape external\+Reactivity} is read both by the spatial kinetics and point kinetics solvers and is used to instantly add (or remove) a certain reactivity at the beginning of a transient. 
\item {\itshape precursor\+Powers} can be read by the point kinetics sub-\/solver in case the user wishes to set initial concentrations of precursors. If not found, precursor concentration are initialized so to be in equilibrium with the starting conditions (i.\+e. a steady state is assumed). 
\end{DoxyItemize}

All Ge\+N-\/\+Foam neutronics models can be used for liquid-\/fuel reactors. One can activate this option using the {\itshape liquid\+Fuel} keyword in $\ast$/system/control\+Dict$\ast$. Of course, in such case one should pay attention to setting proper boundary conditions for the precursors.

A commented reactor\+State can be found in \href{https://gitlab.com/foam-for-nuclear/GeN-Foam/-/tree/master/Tutorials/3D_SmallESFR/rootCase/constant/neutroRegion/reactorState}{\tt 3\+D\+\_\+\+Small\+E\+S\+FR}.

NB\+: Please note that in parallel calculations, the updated {\itshape reactor\+State} can be found in {\itshape processor0/constant/neutro\+Region/}. 

\subsection*{Models}

Neutronics calculations are performed by classes derived from {\itshape neutronics} that contain specific sub-\/solvers\+:
\begin{DoxyItemize}
\item {\itshape point\+Kinetic\+Neutronics} for point kinetics calculations (see {\itshape \hyperlink{pointKineticNeutronics_8H}{point\+Kinetic\+Neutronics.\+H}})
\item {\itshape diffusion\+Neutronics} for diffusion calculations (see {\itshape \hyperlink{diffusionNeutronics_8H}{diffusion\+Neutronics.\+H}})
\item {\itshape adjoint\+Diffusion\+Neutronics} for adjoint diffusion calculations (see {\itshape \hyperlink{adjointDiffusionNeutronics_8H}{adjoint\+Diffusion\+Neutronics.\+H}})
\item {\itshape S\+P3\+Neutronics} for diffusion calculations (see {\itshape \hyperlink{SP3Neutronics_8H}{S\+P3\+Neutronics.\+H}})
\item {\itshape S\+N\+Neutronics} for discrete ordinates calculations (see {\itshape \hyperlink{SNNeutronics_8H}{S\+N\+Neutronics.\+H}}) For the user, the derived classes translate into runtime selectable models. The specific sub-\/solver to be used in a simulation can be selected at runtime in the {\itshape constant/neutro\+Region/neutronics\+Properties} dictionary.
\end{DoxyItemize}

 {\bfseries The {\itshape neutronics\+Properties} dictionary}

The {\itshape neutronics\+Properties} dictionary is found under {\itshape constant/neutro\+Region/} and it can be used to set the type of neutronic simulation by using the following keywords\+: 
\begin{DoxyItemize}
\item {\itshape model} is used to define what type of simulation needs to be performed. It can be {\itshape point\+Kinetics}, {\itshape diffusion\+Neutronics}, {\itshape S\+P3\+Neutronics}, $\ast$$\ast$\+S\+N\+Neutronics, {\itshape adjoint\+Diffusion}. {\itshape adjoint\+Diffusion} has been developed only as an eigenvalue solver. The others can be used for transient calculations. However, the SN transient solver has not been tested. In addition, it is currently not accelrated, thus extremely slow (it can require hundreds of iterations per time step). 
\item {\itshape eigenvalue\+Neutronics} should be set to {\itshape true} for eigenvalue calculations, false for transients. 
\end{DoxyItemize}One can find detailed, commented examples in most tutorials. See for instance \href{https://gitlab.com/foam-for-nuclear/GeN-Foam/-/tree/master/Tutorials/3D_SmallESFR/rootCase/constant/neutroRegion/neutronicsProperties}{\tt 3\+D\+\_\+\+Small\+E\+S\+FR} (single phase). 

\subsection*{Various properties}

In Ge\+N-\/\+Foam, cross-\/sections and several other neutronics properties are handled by the {\itshape \hyperlink{XS_8H}{X\+S.\+H}} class, or by its low-\/memory version {\itshape \hyperlink{XSLowMem_8H}{X\+S\+Low\+Mem.\+H}}.

 {\bfseries The {\itshape nuclear\+Data} dictionaries}

The {\itshape nuclear\+Data} dictionary can be found under {\itshape constant/neutro\+Region/}. It contains all basic nuclear properties for the reference reactor state. The other {\itshape nuclear\+Data...} files in {\itshape constant/neutronics/} should include the cross-\/sections for perturbed reactor states. In addition, these files include information about the perturbed and reference ({\itshape nuclear\+Data}) reactor state. For instance, {\itshape nuclear\+Data\+Fuel\+Temp} must include {\itshape Tfuel\+Ref} and {\itshape Tfuel\+Perturbed}, which represent the temperatures at which the reference ({\itshape nuclear\+Data}) and perturbed ({\itshape nuclear\+Data\+Fuel\+Temp}) cross sections have been calculated, respectively. Linear interpolation is performed by Ge\+N-\/\+Foam between reference and perturbed reactor states, except for fuel temperature, for which a logarithmic or square root interpolation is provided (depending on the spectrum, which in turns is defined by the keyword {\itshape fast\+Neutrons}). If no data are provided, the reference cross sections are used. Nuclear data can be generated using any nuclear code. The \href{https://gitlab.com/foam-for-nuclear/GeN-Foam/-/tree/master/Tools/serpentToFoam/serpent2.1.23}{\tt serpent\+To\+Foam} routines provided with Ge\+N-\/\+Foam (in the {\itshape Tools} folder) is an Octave script that automatically converts Serpent output files into the nuclear data files employed by Ge\+N-\/\+Foam. The entry {\itshape disc\+Factor} is used only if discontinuity factors have to be used. The term {\itshape integral\+Flux}, is used only if the automatic adjustment of discontinuity factors is performed \cite{FIORINA2016212}. Nonetheless, these entries should always be present. ~\newline
~\newline
One can find detailed, commented examples of nuclear\+Data in the tutorials \href{https://gitlab.com/foam-for-nuclear/GeN-Foam/-/tree/master/Tutorials/3D_SmallESFR/rootCase/constant/neutroRegion/nuclearData}{\tt 3\+D\+\_\+\+Small\+E\+S\+FR} (for diffusion or S\+P3), \href{https://gitlab.com/foam-for-nuclear/GeN-Foam/-/tree/master/Tutorials/Godiva_SN/constant/neutroRegion/nuclearData}{\tt Godiva\+\_\+\+SN} (for discrete ordinates) and \href{https://gitlab.com/foam-for-nuclear/GeN-Foam/-/tree/master/Tutorials/2D_onePhaseAndPointKineticsCoupling/rootCase/constant/neutroRegion/nuclearData}{\tt 2\+D\+\_\+one\+Phase\+And\+Point\+Kinetics\+Coupling} (for point kinetics).

N.\+B.\+: cross sections must be expressed according to the International System of Units (so m, not cm).

N.\+B.\+2\+: default\+Prec has 1/m3 units except for the adjoint solver that needs 1/m2/s.

N.\+B.\+3\+: the {\itshape nuclear\+Data...} files must always be present, even when not parametrizing cross-\/sections. If no parametrization is needed, the “zone” card must be left “blank” as\+: {\ttfamily zones();} 

An additional dictionary is needed to provide the quadrature set when performing discrete ordinate calculations.

 {\bfseries The {\itshape quadrature\+Set} dictionary}

The {\itshape quadrature\+Set} dictionary is found under {\itshape constant/neutro\+Region/}. It ccontains the quadrature set for discrete ordinate calculations. ~\newline
~\newline
One can find examples of three different quadrature set in the tutorial \href{https://gitlab.com/foam-for-nuclear/GeN-Foam/-/tree/master/Tutorials/Godiva_SN/constant/neutroRegion/}{\tt Godiva\+\_\+\+SN}. S4 and S8 chebichev Legendre quadrature sets can be found in \href{https://gitlab.com/foam-for-nuclear/GeN-Foam/-/tree/develop/Tools/chebichevLegendreQuadratureSets/}{\tt Godiva\+\_\+\+SN} 

Finally, the {\itshape C\+R\+Move} dictionary can be used to move control rods.

 {\bfseries The {\itshape C\+R\+Move} dictionary}

The C\+R\+Move$\ast$ dictionary can be found under {\itshape constant/neutro\+Region/}. It contains input data for control rods movement. Control rods can be moved from the initial position to a new one by selecting initial and final time of the insertion/extraction and the speed of insertion/extraction (positive speed for insertion). ~\newline
~\newline
One can find a commented example in the tutorial \href{https://gitlab.com/foam-for-nuclear/GeN-Foam/-/tree/master/Tutorials/3D_SmallESFR/rootCase/constant/neutroRegion/CRmove}{\tt 3\+D\+\_\+\+Small\+E\+S\+FR}, though this option is not actually used in the tutorial. 

\subsection*{Initial and boundary conditions}

As in all standard Open\+F\+O\+AM solvers, initial values (IC) and boundary conditions (BC) should be provided in the “0” folder, or in the folder corresponding to the {\itshape start\+Time} of the simulation, if different than 0. In the case of neutronics, the user can either specify different IC and BC for each one of the energy groups (with fluxes that must be named {\itshape flux\+Star0}, {\itshape flux\+Star1}, etc…), or provide the same IC and BC to all fluxes by using the {\itshape default\+Flux} field. In case of S\+P3 calculations, the IC and BC for the second moment can be imposed either for each energy (using fields named {\itshape flux\+Star20}, {\itshape flux\+Star21}, etc…), or to all energies by using the {\itshape default\+Flux2} field. When both {\itshape default\+Flux} and {\itshape flux\+Star...} are present, the solver gives priority to {\itshape flux\+Star...}. In case of SN calculations, it is suggested not to modify the boundary conditions and to use the {\itshape default\+Flux} file (an example is providedin the Godiva\+\_\+\+SN tutorial). When employing the adjoint solver, you will have to add the fields {\itshape adjoint\+Default\+Prec} and {\itshape adjoint\+Default\+Flux} in your initial time.

In addition to the standard Open\+F\+O\+AM BC, an albedo boundary condition (see {\itshape \hyperlink{albedoSP3FvPatchField_8H}{albedo\+S\+P3\+Fv\+Patch\+Field.\+H}}) is available in Ge\+N-\/\+Foam for diffusion and S\+P3 calculations and can be used according to the following syntax\+:


\begin{DoxyCode}
type            albedoSP3;
gamma       0.5; // defined as (1-alpha)/(1+alpha)/2, alpha being the albedo coefficient
diffCoeffName   Dalbedo;  //not to be changed
fluxStarAlbedo  fluxStarAlbedo; //not to be changed
forSecondMoment false;  //true in case it is a condition for a second moment flux (for SP3 calculations)
value           uniform 1;
\end{DoxyCode}
 Please note that the boundary condition needs to be set both for first and second moments in S\+P3.

IC and BC for precursors do not have to be specified for standard reactors. On the other hand, they should be specified in case of liquid fuel reactors (e.\+g., Molten Salt Reactors). This is possible by creating a {\itshape default\+Prec} field, in case the same conditions apply to all precursor groups, or by creating the fields named {\itshape prec0}, {\itshape prec1}, etc., in case different conditions must be provided for different precursor groups.

N.\+B.\+: boundary conditions must be applied to {\itshape flux\+Star...} and not to {\itshape flux...} since Ge\+N-\/\+Foam solves for these variables. {\itshape flux\+Star...} represent continuous fluxes, while {\itshape flux...} represent the real fluxes. They differ only in case discontinuity factors are employed \cite{FIORINA2016212}.

 {\bfseries Setting the weighting in point-\/kinetics calculations}

A correct evaluation of the reactivity worth of delayed neutron precursors in M\+S\+Rs, as well as of the impact of temperatures on reactivities, normally requires the knowledge of the adjoint flux. In Ge\+N-\/\+Foam, the field {\itshape one\+Group\+Flux} is used by the point kinetic solver for weighting temperatures, densities and precursors. When fluxes are not calculated via a spatial neutronics calculation, one has to manually provide the {\itshape one\+Group\+Flux} in {\itshape 0/neutro\+Region}. As an alternative, one can use the {\itshape initial\+One\+Group\+Flux\+By\+Zone} keyword in {\itshape nuclear\+Data} (see \href{https://gitlab.com/foam-for-nuclear/GeN-Foam/-/blob/master/Tutorials/1D_MSR_pointKinetics/rootCase/constant/neutroRegion/nuclearData}{\tt 1\+D\+\_\+\+M\+S\+R\+\_\+point\+Kinetics}). Please notice that\+: 
\begin{DoxyItemize}
\item If calcuclated fluxes are available in {\itshape neutro\+Region}, these will be user to recacluate and overwrite {\itshape one\+Group\+Flux}. 
\item If no fluxes are available, the neutronic-\/subsolver will use the provided {\itshape one\+Group\+Flux} 
\item If the {\itshape initial\+One\+Group\+Flux\+By\+Zone} keyword in used in {\itshape nuclear\+Data}, this will be used to overwrite {\itshape one\+Group\+Flux} 
\end{DoxyItemize}

\subsection*{Discretization and solution}

Details for discretization and solution of equations are handled in a standard Open\+F\+O\+AM way, i.\+e., through the {\itshape fv\+Solution} and {\itshape fv\+Schemes} dictionaries in {\itshape constant/neutro\+Region}.

© All rights reserved. E\+C\+O\+LE P\+O\+L\+Y\+T\+E\+C\+H\+N\+I\+Q\+UE F\+E\+D\+E\+R\+A\+LE DE L\+A\+U\+S\+A\+N\+NE, Switzerland, 2021 